\documentclass[a4paper,11pt]{article}

\usepackage[T1]{fontenc}
\usepackage{amsmath,amssymb}
\usepackage{booktabs}
\usepackage[margin=1in]{geometry}
\usepackage{xcolor}
\usepackage{listings}
\usepackage[
    colorlinks=true,
    citecolor=blue,
    linkcolor=blue,
    urlcolor=magenta
]{hyperref}

\newcommand{\matlab}{\textsf{MATLAB}}
\newcommand{\python}{\textsf{Python}}
\newcommand{\file}[1]{\texttt{\detokenize{#1}}}

\lstdefinestyle{code}{
    basicstyle=\ttfamily\small,
    keywordstyle=\color{blue!70!black},
    commentstyle=\color{green!45!black},
    stringstyle=\color{purple!70!black},
    columns=fullflexible,
    keepspaces=true,
    breaklines=true,
    showstringspaces=false,
    frame=single,
    framerule=0.3pt,
    xleftmargin=0.5em,
    xrightmargin=0.5em
}

\title{\matlab\ and \python\ Files for Pareto Extrapolation}
\author{Alexis Akira Toda\thanks{Department of Economics, Emory University.
    Email: \href{mailto:alexis.akira.toda@emory.edu}{alexis.akira.toda@emory.edu}.}}
\date{Revised August 16, 2026}

\begin{document}

\maketitle

\section{Introduction}

This package implements the Pareto extrapolation algorithm developed by
\cite{GouinBonenfantToda2023}.  The method augments a conventional finite-grid
approximation with an analytical Pareto tail.  It computes the Pareto exponent,
the stationary joint distribution of an exogenous state and a size variable
such as wealth, and top wealth shares.  The characterization of the Pareto
exponent is based on \cite{BeareToda2022}.

The source repository is
\href{https://github.com/alexisakira/Pareto-extrapolation}{github.com/alexisakira/Pareto-extrapolation}.
The software is maintained by Alexis Akira Toda and is provided without
warranty.  Users should verify numerical accuracy and economic assumptions for
their own application.  If these routines are used in research, please cite
\cite{GouinBonenfantToda2023}.

The folder contains parallel implementations:

\begin{lstlisting}[style=code]
PE/
|-- matlab/                  MATLAB implementation and example
|-- python/
|   |-- src/pareto_extrapolation/   Python package
|   |-- examples/                   Python example
|   |-- tests/                      MATLAB-parity and edge-case tests
|   `-- pyproject.toml
|-- readme.tex
`-- readme.pdf
\end{lstlisting}

\subsection*{Index and probability conventions}

Let $s,t\in\{1,\ldots,S\}$ denote the current and next exogenous states,
$j\in\{1,\ldots,J\}$ a transitory shock, and $n\in\{1,\ldots,N\}$ a grid
index.  Whenever an $S^2\times J$ array is required, the row for $(s,t)$ is

\begin{equation*}
    r=S(s-1)+t
\end{equation*}

in one-based notation.  Thus the rows are ordered as

\begin{equation*}
    (1,1),\ldots,(1,S);\quad\ldots;\quad(S,1),\ldots,(S,S).
\end{equation*}

In the joint transition matrix, the grid index varies fastest:

\begin{equation*}
    (s,n)=(1,1),\ldots,(1,N);\ldots;(S,1),\ldots,(S,N).
\end{equation*}

Probability transition matrices are row-stochastic.  The survival probability
$V(s,t)$ is conditional on the transition from $s$ to $t$; use $V=1$ for an
infinitely lived model.  Numerical inputs must be real and finite.

\section{MATLAB}

The \matlab\ implementation is in the \file{matlab} folder.  It uses standard
\matlab\ functionality, including sparse matrices, \file{eigs},
\file{discretize}, \file{pchip}, and \file{interp1}, and has been tested with
\matlab\ R2024b.  From the \file{PE} folder, run:

\begin{lstlisting}[style=code,language=Matlab]
addpath('matlab')
example
\end{lstlisting}

\subsection{Exponential grid: \file{expGrid.m}}

\begin{lstlisting}[style=code,language=Matlab]
grid = expGrid(a,b,c,N)
\end{lstlisting}

The function returns a strictly increasing $1\times N$ shifted-log grid on the
closed interval $[a,b]$, with \file{grid(1) = a} and
\file{grid(end) = b} exactly.  The concentration parameter must satisfy

\begin{equation*}
    a<c<\frac{a+b}{2},
\end{equation*}

and $N$ must be an integer no smaller than two.  For odd $N$, $c$ is the middle
grid point; for even $N$, it lies between the middle points in shifted-log
space.  This differs from releases through 2021, which excluded $a$ and
generated points on $(a,b]$.

\subsection{Pareto exponent: \file{getZeta.m}}

\begin{lstlisting}[style=code,language=Matlab]
[zeta,typeDist] = getZeta(PS,PJ,V,G,zetaBound)
\end{lstlisting}

The inputs are:

\begin{itemize}
    \item \file{PS}: the $S\times S$ transition matrix of the exogenous state;
    \item \file{PJ}: conditional shock probabilities with size $1\times J$,
    $S\times J$, or $S^2\times J$.  The shorter forms assume independence
    from $(s,t)$ or dependence only on $s$;
    \item \file{V}: a scalar or $S\times S$ matrix of survival probabilities;
    \item \file{G}: asymptotic gross growth rates with size $S\times J$ or
    $S^2\times J$; and
    \item \file{zetaBound}: an optional two-element search interval, with
    default $[10^{-2},100]$.
\end{itemize}

Define

\begin{align*}
    M_{st}(z) &= \sum_{j=1}^{J}\Pr(j\mid s,t)G_{stj}^{z},\\
    A(z) &= PS\odot V\odot M(z),
\end{align*}

where $\odot$ denotes elementwise multiplication.  The Pareto exponent solves

\begin{equation*}
    \rho\!\left(A(\zeta)\right)=1,
\end{equation*}

and \file{typeDist} is the normalized left Perron vector of $A(\zeta)$.  A
root exactly on a search bound is accepted.  If no root exists inside the
interval, the function returns the relevant endpoint, issues a warning, and
returns \file{typeDist = []}.

\subsection{Joint transition matrix: \file{getQ.m}}

\begin{lstlisting}[style=code,language=Matlab]
[Q,piStar] = getQ(PS,PJ,V,x0,xGrid,gstjn,Gstj,zeta)
\end{lstlisting}

The first three inputs follow \file{getZeta.m}.  The remaining inputs are:

\begin{itemize}
    \item \file{x0}: the size of a newborn or reset agent.  It is used when at
    least one entry of $V$ is below one and must satisfy $x_1\leq x_0<x_N$;
    \item \file{xGrid}: a strictly increasing $N$-element grid with $x_N>0$;
    \item \file{gstjn}: the law of motion evaluated on the grid.  It has $JN$
    columns, with columns $N(j-1)+1$ through $Nj$ corresponding to shock $j$.
    It may have $S$ rows or $S^2$ rows;
    \item \file{Gstj}: optional asymptotic slopes with $S$ or $S^2$ rows.
    Pass \file{[]} to estimate them from the two largest grid points; and
    \item \file{zeta}: the optional Pareto exponent.  Pass \file{[]} to call
    \file{getZeta.m} internally.
\end{itemize}

The function returns the $SN\times SN$ transition matrix \file{Q} as a sparse
matrix.  If requested, \file{piStar} is the full $SN\times1$ stationary
distribution.  The marginal distribution on the size grid is

\begin{lstlisting}[style=code,language=Matlab]
xDist = sum(reshape(piStar,N,S),2);
\end{lstlisting}

The function checks the row sums of \file{Q} and the stationary-distribution
residual.  Pareto extrapolation requires $\zeta>1$ for a finite mean.  If
$\zeta\leq1$, the function warns and returns NaN-valued outputs.

\subsection{Top wealth shares: \file{getTopShares.m}}

\begin{lstlisting}[style=code,language=Matlab]
topShare = getTopShares(topProb,wGrid,wDist,zeta)
\end{lstlisting}

\file{topProb} is a strictly increasing vector in $[0,1]$,
\file{wGrid} is a strictly increasing wealth grid, and \file{wDist} is a
probability vector on that grid.  Omit \file{zeta}, pass \file{[]}, or pass
zero to use upper-grid truncation; otherwise $\zeta$ must exceed one.
The output has the same row or column orientation as \file{topProb}.

Interpolation uses the shape-preserving \file{pchip} method.  Under Pareto
extrapolation, the wealth represented by the largest grid point is valued as

\begin{equation*}
    \widetilde w_N=\frac{\zeta}{\zeta-1}w_N.
\end{equation*}

Negative wealth is allowed, but aggregate net wealth---including the Pareto
correction when applicable---must be strictly positive.  Top shares can
legitimately exceed one when households outside the selected group have debt,
and the top-share curve can decline as indebted households enter the group.
The function therefore does not clamp the result to $[0,1]$ or force
monotonicity.

\subsection{Example: \file{example.m}}

The example uses a two-state return process and a two-point transitory shock:

\begin{lstlisting}[style=code,language=Matlab]
xGrid = expGrid(xMin,xMax,x0,N);
gstjn = kron(Gstj,xGrid);

[zeta,typeDist] = getZeta(PS,PJ,V,Gstj,[0.1 10]);
[Q,piStar] = getQ(PS,PJ,V,x0,xGrid,gstjn,Gstj,zeta);

xDist = sum(reshape(piStar,N,S),2);
topShare = getTopShares([0.001 0.01 0.1],xGrid,xDist,zeta);
\end{lstlisting}

It produces $\zeta\approx1.721899$ and top $0.1\%$, $1\%$, and $10\%$ shares
of approximately $0.073619$, $0.177085$, and $0.426102$.  The script also
sets session-wide default graphics formatting and plots the wealth-tail
probability.

\section{Python}

The \python\ implementation is an installable NumPy/SciPy package in the
\file{python} folder.  It requires \python\ 3.10 or later.  From that folder,
install the package and optional example and test dependencies with

\begin{lstlisting}[style=code]
python -m pip install -e ".[example,test]"
pytest
python examples/example.py
\end{lstlisting}

The public functions use snake-case names and are imported from the package as
follows:

\begin{lstlisting}[style=code,language=Python]
from pareto_extrapolation import (
    exp_grid,
    get_q,
    get_top_shares,
    get_zeta,
)
\end{lstlisting}

The state-pair and joint-state orderings are identical to the \matlab\ version.
Python arrays are zero-indexed, but users do not need to transform the data
ordering.  One-dimensional arrays replace \matlab\ row or column vectors.

\subsection{Exponential grid: \file{exp_grid}}

\begin{lstlisting}[style=code,language=Python]
grid = exp_grid(a, b, c, n)
\end{lstlisting}

This is the direct NumPy counterpart of \file{expGrid.m}.  It returns a
one-dimensional array containing both endpoints exactly and applies the same
input restrictions and shifted-log formula.

\subsection{Pareto exponent: \file{get_zeta}}

\begin{lstlisting}[style=code,language=Python]
zeta, type_dist = get_zeta(ps, pj, v, g, zeta_bound=(1e-2, 100.0))
\end{lstlisting}

The array shapes and probability conventions match \file{getZeta.m}.  A
one-dimensional \file{pj} is interpreted as a state-independent shock
distribution.  The implementation uses dense NumPy eigenvalue calculations
for the small $S\times S$ matrix and SciPy's bracketed root solver.  It issues
\file{RuntimeWarning} and returns an empty NumPy array when the exponent lies
outside the search interval.

\subsection{Joint transition matrix: \file{get_q}}

\begin{lstlisting}[style=code,language=Python]
q, pi_star = get_q(
    ps, pj, v, x0, x_grid, gstjn,
    g_stj=None, zeta=None,
    compute_stationary=True,
)
\end{lstlisting}

The input meanings and array order match \file{getQ.m}.  The function returns
\file{q} as a SciPy CSR sparse matrix and \file{pi_star} as a one-dimensional
NumPy array.  To skip the stationary calculation, use

\begin{lstlisting}[style=code,language=Python]
q, pi_star = get_q(..., compute_stationary=False)
\end{lstlisting}

The second returned value is then \file{None}.  If \file{g_stj} or
\file{zeta} is \file{None}, it is estimated or computed internally.

The marginal grid distribution is

\begin{lstlisting}[style=code,language=Python]
x_dist = pi_star.reshape(S, N).sum(axis=0)
\end{lstlisting}

Unlike the legacy-compatible \matlab\ interface, the \python\ function raises
\file{ValueError} when $\zeta\leq1$ instead of returning a NaN matrix.

\subsection{Top wealth shares: \file{get_top_shares}}

\begin{lstlisting}[style=code,language=Python]
top_share = get_top_shares(top_prob, w_grid, w_dist, zeta=0.0)
\end{lstlisting}

This function follows the revised negative-wealth behavior of
\file{getTopShares.m} and uses SciPy's shape-preserving PCHIP interpolator.
A scalar probability input returns a float.  A one-dimensional input returns a
NumPy array.  The \matlab\ and SciPy PCHIP implementations can differ at about
the $10^{-10}$ level even when their interpolation knots agree.

\subsection{Example and tests}

The script \file{python/examples/example.py} reproduces the \matlab\ example,
prints the exponent, type distribution, top shares, and sparse-matrix size, and
plots the wealth tail with Matplotlib.

The tests in \file{python/tests} cover endpoint construction, exact-bound
roots, invalid inputs, negative wealth, zero mass at the largest grid point,
the complete economic example, and a fixed transition matrix generated by the
\matlab\ implementation.  The transition entries agree to floating-point
precision; the end-to-end top shares agree within $10^{-10}$.

\begin{thebibliography}{9}
\small
\raggedright

\bibitem{BeareToda2022}
Brendan K. Beare and Alexis Akira Toda.
``Determination of Pareto Exponents in Economic Models Driven by Markov
Multiplicative Processes.''
\emph{Econometrica}, 90(4):1811--1833, 2022.
\href{https://doi.org/10.3982/ECTA17984}{doi:10.3982/ECTA17984}.

\bibitem{GouinBonenfantToda2023}
{\'E}milien Gouin-Bonenfant and Alexis Akira Toda.
``Pareto Extrapolation: An Analytical Framework for Studying Tail Inequality.''
\emph{Quantitative Economics}, 14(1):201--233, 2023.
\href{https://doi.org/10.3982/QE1817}{doi:10.3982/QE1817}.

\end{thebibliography}

\end{document}
